\documentclass[12pt,a4paper]{article}
\usepackage[utf8]{inputenc}
\usepackage[T5]{fontenc}
\usepackage{geometry}
\usepackage{setspace}
\usepackage{graphicx}
\usepackage{booktabs}
\usepackage{float}
\usepackage{amsmath}
\usepackage{tabularx}
\usepackage{longtable}
\usepackage{array}
\usepackage{hyperref}
\geometry{margin=1in}
\onehalfspacing
\setcounter{tocdepth}{2}
\renewcommand{\arraystretch}{1.15}

\title{Báo cáo Dự án Xử lý Ngôn ngữ Tự nhiên\\
Phiên bản tháng 03/2026}
\author{Team SkyDom+}
\date{Tháng 3 năm 2026}

\begin{document}

\begin{titlepage}
\centering
\vspace*{1.0cm}
{\LARGE Phân tích cảm xúc và trích xuất vấn đề từ đánh giá thương mại điện tử\par}
\vspace{0.2cm}
{\LARGE bằng NLP cổ điển và mở rộng Transformer\par}
\vspace{0.8cm}
{\Large Báo cáo học thuật dự án\par}
\vspace{0.8cm}
{\large Đội: Team SkyDom+\par}
\vspace{0.3cm}
{\large Thành viên nhóm:\par}
\vspace{0.4cm}
{\large
\begin{tabular}{@{}>{\raggedright\arraybackslash}p{5.8cm} >{\raggedright\arraybackslash}p{2.7cm}@{}}
Nguyen Nhu Tien Tuan & 22BA13314 \\
Tran Hien Chuong & 22BA13056 \\
Do Gia Bao & 22BA13042 \\
Bui The Anh & 22BA13007 \\
Le Duc An & 22BA13002 \\
Pham Quang Tung & 22BA13320 \\
Khuong Bao Duc & 22BA13077 \\
\end{tabular}\par}
\vfill
{\large Ngày cập nhật: 12/03/2026\par}
\end{titlepage}

\pagenumbering{roman}

\begin{abstract}
Báo cáo này trình bày một nghiên cứu NLP theo cấu trúc học thuật thống nhất. Hệ thống gồm ba nhiệm vụ: phân loại cảm xúc nhị phân có cơ chế bất định, trích xuất issue đa nhãn, và mở rộng bằng transformer. Dữ liệu nguồn gồm 152.410 review; bộ dữ liệu issue đã gán nhãn gồm 19.000 dòng. Với sentiment cổ điển, biến thể V6 (word 1--2 + char 3--5 + negation, class weight w10) đạt test \texttt{precision\_0=0.677}, \texttt{recall\_0=0.932}, \texttt{F2\_0=0.867}, coverage 0.885 ở dải ngưỡng 0.40/0.60. Mô hình transformer sentiment đạt \texttt{accuracy=0.966}, \texttt{precision\_0=0.799}, \texttt{recall\_0=0.853}, \texttt{F2\_0=0.842}, coverage 0.910. Với issue extraction trên tập test 3.800 mẫu, OVR Logistic đạt \texttt{micro-F1=0.780}, \texttt{macro-F1=0.387}; OVR LinearSVM đạt \texttt{micro-F1=0.771}, \texttt{macro-F1=0.405}; trong khi dummy baseline có micro-F1 cao (\texttt{0.809}) nhưng macro-F1 rất thấp (\texttt{0.101}) do lệch nhãn \texttt{other}. Gói nâng cấp không thêm dữ liệu mở rộng thêm benchmark giáo trình, n-gram LM, bootstrap/significance, và thử nghiệm transformer issue + hybrid routing. Toàn bộ số liệu trong báo cáo lấy trực tiếp từ artifact thực nghiệm trong thư mục \texttt{results/}.
\end{abstract}

\tableofcontents
\newpage
\pagenumbering{arabic}

\section{Giới thiệu}
Ngôn ngữ trong review ngắn, nhiễu và thường trộn nhiều sắc thái trong cùng một câu. Vì vậy bài toán được triển khai như một pipeline NLP hoàn chỉnh thay vì chỉ so sánh một mô hình phân loại đơn lẻ \cite{jurafsky2026,turney2002}.

\subsection{Mục tiêu dự án}
Mục tiêu thứ nhất là xây dựng pipeline phân tích cảm xúc có thể phân biệt negative và positive, đồng thời xử lý các trường hợp uncertain cẩn thận hơn so với bộ phân loại gán nhãn cứng thông thường. Mục tiêu thứ hai là đi xa hơn nhãn cảm xúc tổng quát để trích xuất các nhóm issue cụ thể phục vụ xử lý tiếp theo. Mục tiêu thứ ba là đối chiếu hướng NLP sparse cổ điển với hướng transformer theo ngữ cảnh sâu trên cùng một miền dữ liệu review.

\subsection{Đối tượng sử dụng}
Nhóm người dùng chính là bộ phận chăm sóc khách hàng, vận hành, và phân tích sản phẩm khi phải xử lý khối lượng lớn phản hồi thương mại điện tử. Nhóm đối tượng thứ hai là giảng viên và sinh viên trong học phần NLP, nơi dự án có thể đóng vai trò tình huống minh họa cho phân loại văn bản, xử lý bất định, học đa nhãn và mở rộng transformer trong cùng một bài toán ứng dụng.

Ở góc độ vận hành thực tế, bài toán này có các ``nỗi đau'' rõ ràng. Thứ nhất, việc đọc thủ công số lượng review lớn tốn thời gian, khiến các phản hồi tiêu cực khẩn cấp dễ bị chìm trong số lượng lớn review tích cực. Thứ hai, nhãn cảm xúc tổng quát là chưa đủ cho hành động, vì bộ phận hỗ trợ cần biết cụ thể vấn đề thuộc nhóm giao hàng, kích hoạt, hoàn tiền hay nghi ngờ gian lận. Thứ ba, review quá ngắn hoặc mơ hồ làm tăng rủi ro quyết định; nếu ép gán nhãn cứng, hệ thống có thể vừa tăng cảnh báo sai vừa bỏ sót khiếu nại thật.

Các ràng buộc này dẫn tới cách tiếp cận vượt ra ngoài việc tối ưu một chỉ số accuracy đơn lẻ. Báo cáo ưu tiên đánh giá nhạy theo recall cho lớp negative, trích xuất issue đa nhãn để phục vụ xử lý tiếp theo, và cơ chế uncertain cho trường hợp thiếu thông tin. Nhờ đó, lựa chọn mô hình NLP gắn trực tiếp với chất lượng phản hồi dịch vụ và thời gian xử lý.

Phạm vi nâng cấp bám nguyên tắc không thêm dữ liệu mới: chỉ cải tiến xử lý, mô hình hóa, hiệu chỉnh và đánh giá.

\section{Định nghĩa bài toán NLP}
\subsection{Task A: Hiểu cảm xúc}
Task A là phân loại cảm xúc nhị phân từ văn bản review. Rating 1--2 ánh xạ về lớp 0 (negative), rating 4--5 ánh xạ về lớp 1 (positive), rating 3 dành cho phân tích bất định. Do dữ liệu mất cân bằng mạnh (xấp xỉ 9,5:1), tiêu chí chọn mô hình theo thứ tự ưu tiên: \texttt{recall\_0} rồi \texttt{F2\_0} rồi \texttt{precision\_0}.

\subsection{Task B: Trích xuất issue}
Task B là gán đa nhãn issue với 9 nhãn (delivery, redemption, product quality, customer service, refund, usability, value/price, fraud/scam, other). Đây là ánh xạ từ văn bản không cấu trúc sang tín hiệu vận hành có cấu trúc \cite{tsoumakas2011,zhang2014}.

\subsection{Task C: Mở rộng ngữ cảnh sâu}
Task C dùng transformer fine-tuning để kiểm tra năng lực mô hình hóa ngữ cảnh sâu, xử lý case khó và so sánh trực tiếp với hướng sparse cổ điển \cite{vaswani2017,devlin2019}.

\section{Dữ liệu và đặc trưng ngôn ngữ}
\subsection{Nguồn dữ liệu và quy mô}
Tập nguồn có 152.410 review. Phân bố rating được trình bày ở Bảng \ref{tab:rating_distribution_vi}.

\begin{table}[H]
\centering
\caption{Phân bố rating trong dữ liệu nguồn}
\label{tab:rating_distribution_vi}
\begin{tabular}{lr}
\toprule
Rating & Số lượng \\
\midrule
1 & 12.326 \\
2 & 1.873 \\
3 & 3.271 \\
4 & 6.692 \\
5 & 128.248 \\
\bottomrule
\end{tabular}
\end{table}

Phần strong-label sentiment (rating 1,2,4,5) có 149.139 dòng, gồm 14.199 negative và 134.940 positive.

\subsection{Nhiễu ngôn ngữ quan sát được}
Dữ liệu có nhiều dạng nhiễu: viết tắt (\textit{gr8}, \textit{thx}, \textit{idk}), câu rất ngắn, lỗi chính tả, và câu trộn polarity như \textit{good but late delivery}. Các hiện tượng negation và contrast xuất hiện thường xuyên.

\subsection{Dữ liệu gán nhãn issue}
Bộ issue hợp nhất có 19.000 dòng. Tần suất nhãn được trình bày ở Bảng \ref{tab:issue_freq_vi}.

\begin{table}[H]
\centering
\caption{Tần suất nhãn issue trong bộ gán nhãn}
\label{tab:issue_freq_vi}
\begin{tabular}{lrr}
\toprule
Nhãn & Số lượng & Tỷ lệ \\
\midrule
delivery\_shipping & 972 & 0,0512 \\
redemption\_activation & 284 & 0,0149 \\
product\_quality & 1.251 & 0,0658 \\
customer\_service & 182 & 0,0096 \\
refund\_return & 93 & 0,0049 \\
usability & 873 & 0,0459 \\
value\_price & 558 & 0,0294 \\
fraud\_scam & 144 & 0,0076 \\
other & 15.960 & 0,8400 \\
\bottomrule
\end{tabular}
\end{table}

\section{Thiết kế xử lý văn bản}
\subsection{Chuẩn hóa}
Pipeline chuẩn hóa gồm chuyển chữ thường, loại URL, chuẩn hóa dấu câu và khoảng trắng. Mục tiêu là giảm sparsity nhưng vẫn giữ tín hiệu cảm xúc.

\subsection{Xử lý phủ định}
Pipeline sentiment và issue hỗ trợ truyền dấu phủ định theo cửa sổ cục bộ (cấu hình issue hiện tại: cửa sổ 3 token). Đây là thành phần quan trọng để xử lý lật polarity.

\subsection{Xử lý tương phản}
Các tín hiệu tương phản như \textit{but}, \textit{however} được hấp thụ thông qua n-gram biểu diễn cụm từ, giúp mô hình học chuyển pha sắc thái trong câu.

\subsection{Mô hình hóa mức ký tự}
Character n-gram (3--5) được kết hợp với word n-gram để tăng độ bền trước lỗi chính tả, viết tắt và biến thể bề mặt.

\section{Thiết kế đặc trưng}
\subsection{Biểu diễn sparse cổ điển}
Biểu diễn chính cho sentiment là TF-IDF word n-gram (1--2) + char n-gram (3--5), sau đó chọn đặc trưng bằng Chi-square. Biến thể tốt nhất V6 dùng \texttt{k=10000}.

Với issue extraction, pipeline dùng TF-IDF, hỗ trợ char n-gram tùy chọn, và Chi-square với \texttt{k=2000} trước tầng OVR.

\subsection{Khối đặc trưng lexicon}
Trong trạng thái hiện tại, mô hình tốt nhất là sparse data-driven. Các thành phần nghiêng về lexicon được giữ như khối phân tích bổ trợ trong gói mở rộng, không phải đường triển khai chính.

\section{Hướng mô hình hóa}
\subsection{Bộ phân loại sentiment}
Backbone sentiment cổ điển là Logistic Regression. Class-weight được dò trên \texttt{none}, \texttt{balanced}, \texttt{w2}, \texttt{w5}, \texttt{w10}. Các biến thể V0--V7 phản ánh bật/tắt xử lý và đặc trưng.

\subsection{Mô hình issue extraction}
Issue extraction dùng OVR Logistic Regression có hiệu chỉnh xác suất (sigmoid) và tuning ngưỡng theo từng nhãn. OVR LinearSVM và mô hình blend xác suất LR+SVM (\texttt{ovr\_blend\_lr\_svm}) là nhóm so sánh mạnh, đồng thời có thêm hai dummy baseline (\texttt{dummy\_labelset\_majority}, \texttt{dummy\_label\_prior}) để kiểm tra sanity.

\subsection{Mở rộng Transformer}
Khối transformer dùng HuggingFace fine-tuning (họ DistilBERT trong cấu hình này) và đánh giá cùng chính sách ngưỡng bất định để so sánh trực tiếp với mô hình cổ điển.

\section{Thiết lập thực nghiệm}
\subsection{Chia tập và giao thức}
Sentiment huấn luyện trên tập strong-label (rating 1,2,4,5), giữ validation và test độc lập; rating 3 chỉ dùng cho phân tích bất định. Issue extraction chia cố định 13.300/1.900/3.800 (train/val/test). Khối issue-transformer báo cáo trên 3.000/1.000/2.000.

\subsection{Chỉ số chính}
Sentiment dùng các chỉ số \texttt{accuracy}, \texttt{F1}, \texttt{precision\_0}, \texttt{recall\_0}, \texttt{F1\_0}, \texttt{F2\_0}. Issue extraction dùng micro-F1, macro-F1, subset accuracy, hamming loss và F1 theo từng nhãn.

\subsection{Quy tắc bất định}
Selective classification dùng dải ngưỡng \([t_{low}, t_{high}]\). Mẫu nằm trong dải được gắn uncertain; ngoài dải là covered decisions. Báo cáo gồm coverage, selective precision\_0, selective recall\_0, selective F2\_0 và tỷ lệ FN trên phần covered \cite{geifman2017}.

\section{Gói nâng cấp không thêm dữ liệu}
Các thành phần nâng cấp chính trong gói thực nghiệm hiện tại:
\begin{enumerate}
\item Level-3 issue: dò class-weight theo từng nhãn, calibration, và kiểm tra độ ổn định theo nhiễu ngưỡng.
\item Level-4 issue transformer và hybrid routing có giới hạn tỷ lệ route.
\item Bộ benchmark giáo trình: Naive Bayes, Perceptron, LinearSVM, SGD log-loss, SVD semantic baseline, FFNN MLP.
\item Module n-gram language model với add-k smoothing và perplexity.
\item Gói đánh giá nghiêm ngặt: bootstrap CI, kiểm định cặp mô hình, taxonomy lỗi.
\item Ma trận bao phủ học phần sinh trực tiếp từ artifact hiện có.
\end{enumerate}

\section{Kết quả cho sentiment NLP cổ điển}
\subsection{Chọn baseline theo K}
Bảng \ref{tab:chi2_k_vi} cho thấy cấu hình tốt nhất theo từng mức Chi-square K dưới tiêu chí negative-first.

\begin{table}[H]
\centering
\caption{Baseline sentiment tốt nhất theo Chi-square K (validation)}
\label{tab:chi2_k_vi}
\begin{tabular}{lccccc}
\toprule
K & Class weight & Precision\_0 & Recall\_0 & F2\_0 & Accuracy \\
\midrule
1.000 & w10 & 0,555 & 0,910 & 0,806 & 0,922 \\
2.000 & w10 & 0,584 & 0,920 & 0,825 & 0,930 \\
5.000 & balanced & 0,606 & 0,935 & 0,843 & 0,936 \\
10.000 & balanced & 0,621 & 0,936 & 0,850 & 0,940 \\
\bottomrule
\end{tabular}
\end{table}

\subsection{So sánh biến thể}
Bảng \ref{tab:classic_variants_vi} báo cáo kết quả test các biến thể V0--V7; V6 là mô hình chính hiện tại.

\begin{table}[H]
\centering
\caption{So sánh biến thể sentiment cổ điển trên test}
\label{tab:classic_variants_vi}
\begin{tabular}{lccccc}
\toprule
Biến thể & K & Class weight & Precision\_0 & Recall\_0 & F2\_0 \\
\midrule
V0 & 10.000 & balanced & 0,638 & 0,937 & 0,856 \\
V1 & 10.000 & balanced & 0,643 & 0,932 & 0,855 \\
V2 & 5.000 & w10 & 0,647 & 0,926 & 0,852 \\
V3 & 10.000 & balanced & 0,626 & 0,923 & 0,843 \\
V4 & 10.000 & balanced & 0,635 & 0,921 & 0,845 \\
V5 & 10.000 & balanced & 0,653 & 0,933 & 0,859 \\
V6 & 10.000 & w10 & 0,677 & 0,932 & 0,867 \\
V7 & 5.000 & w10 & 0,657 & 0,927 & 0,857 \\
\bottomrule
\end{tabular}
\end{table}

\subsection{Hành vi theo ngưỡng và bất định}
Với V6, dải ngưỡng khuyến nghị là 0,40/0,60. Bảng \ref{tab:classic_threshold_vi} thể hiện trade-off theo dịch chuyển ngưỡng.

\begin{table}[H]
\centering
\caption{Quét ngưỡng của sentiment cổ điển trên test}
\label{tab:classic_threshold_vi}
\begin{tabular}{lcccc}
\toprule
Ngưỡng & Coverage & Sel. precision\_0 & Sel. recall\_0 & Sel. F2\_0 \\
\midrule
0,40 / 0,60 & 0,885 & 0,730 & 0,956 & 0,900 \\
0,45 / 0,60 & 0,891 & 0,708 & 0,956 & 0,894 \\
0,50 / 0,65 & 0,886 & 0,681 & 0,963 & 0,889 \\
0,55 / 0,70 & 0,881 & 0,646 & 0,970 & 0,882 \\
\bottomrule
\end{tabular}
\end{table}

\section{Kết quả cho mở rộng Transformer NLP}
\subsection{Chỉ số Transformer}
Bảng \ref{tab:tf_metrics_vi} tóm tắt kết quả sentiment transformer.

\begin{table}[H]
\centering
\caption{Hiệu năng sentiment transformer}
\label{tab:tf_metrics_vi}
\begin{tabular}{lcccccc}
\toprule
Tập & Accuracy & F1 & Precision\_0 & Recall\_0 & F2\_0 & Coverage \\
\midrule
Validation & 0,966 & 0,981 & 0,793 & 0,870 & 0,854 & 0,906 \\
Test & 0,966 & 0,981 & 0,799 & 0,853 & 0,842 & 0,910 \\
\bottomrule
\end{tabular}
\end{table}

\subsection{Quét ngưỡng Transformer}
\begin{table}[H]
\centering
\caption{Quét ngưỡng transformer trên test}
\label{tab:tf_threshold_vi}
\begin{tabular}{lcccc}
\toprule
Ngưỡng & Coverage & Sel. precision\_0 & Sel. recall\_0 & Sel. F2\_0 \\
\midrule
0,40 / 0,60 & 0,910 & 0,831 & 0,846 & 0,843 \\
0,45 / 0,60 & 0,911 & 0,827 & 0,847 & 0,843 \\
0,50 / 0,65 & 0,912 & 0,821 & 0,850 & 0,844 \\
0,55 / 0,70 & 0,911 & 0,818 & 0,855 & 0,847 \\
\bottomrule
\end{tabular}
\end{table}

\subsection{So sánh hard-case}
Kết quả hard-case cho thấy tính bổ sung giữa hai hướng. Với \textit{not bad} và \textit{good but late delivery}, mô hình cổ điển dự đoán negative còn transformer dự đoán positive. Với chuỗi rất ngắn (\textit{gr8}, \textit{thx}, \textit{idk}), cả hai đều đẩy vào uncertain theo chính sách short-text.

\section{Kết quả cho NLP Issue Extraction}
\subsection{Kết quả tổng thể đa nhãn}
Bảng \ref{tab:issue_overall_vi} trình bày kết quả pipeline issue chính hiện tại (test 3.800 mẫu).

\begin{table}[H]
\centering
\caption{Chỉ số tổng thể issue extraction}
\label{tab:issue_overall_vi}
\begin{tabular}{lcccc}
\toprule
Mô hình & Micro-F1 & Macro-F1 & Subset accuracy & Hamming loss \\
\midrule
OVR Logistic (test) & 0,780 & 0,387 & 0,634 & 0,0585 \\
OVR LinearSVM (test) & 0,771 & 0,405 & 0,617 & 0,0623 \\
OVR Blend LR+SVM (test) & 0,768 & 0,391 & 0,626 & 0,0620 \\
Dummy labelset majority (test) & 0,809 & 0,101 & 0,811 & 0,0438 \\
Dummy label prior (test) & 0,809 & 0,101 & 0,811 & 0,0438 \\
\bottomrule
\end{tabular}
\end{table}

Dummy baseline có thể cao ở micro-F1 và subset accuracy vì gần như luôn dự đoán nhãn \texttt{other}; macro-F1 cho thấy chúng không xử lý tốt các nhãn thiểu số.

\subsection{Chênh lệch theo từng nhãn}
Bảng \ref{tab:issue_delta_vi} so sánh F1 từng nhãn giữa SVM và Logistic trên test.

\begin{table}[H]
\centering
\caption{Chênh lệch F1 theo nhãn (SVM trừ Logistic)}
\label{tab:issue_delta_vi}
\begin{tabular}{lccc}
\toprule
Nhãn & Logistic F1 & SVM F1 & $\Delta$ \\
\midrule
delivery\_shipping & 0,442 & 0,446 & +0,004 \\
redemption\_activation & 0,298 & 0,287 & -0,010 \\
product\_quality & 0,307 & 0,282 & -0,025 \\
customer\_service & 0,225 & 0,168 & -0,057 \\
refund\_return & 0,435 & 0,387 & -0,048 \\
usability & 0,459 & 0,462 & +0,003 \\
value\_price & 0,258 & 0,313 & +0,056 \\
fraud\_scam & 0,150 & 0,391 & +0,241 \\
other & 0,911 & 0,911 & +0,001 \\
\bottomrule
\end{tabular}
\end{table}

\subsection{Chiến lược ngưỡng trong bối cảnh đa nhãn}
Ngưỡng theo từng nhãn của pipeline chính \texttt{ovr\_logreg} được trình bày trong Bảng \ref{tab:issue_thresholds_vi}. Các nhãn hiếm được đặt ngưỡng thấp hơn để tăng recall; \texttt{fraud\_scam} và \texttt{other} được giữ ngưỡng thận trọng hơn.

\begin{table}[H]
\centering
\caption{Ngưỡng cuối cùng theo từng nhãn của pipeline issue}
\label{tab:issue_thresholds_vi}
\begin{tabular}{lc}
\toprule
Nhãn & Ngưỡng \\
\midrule
delivery\_shipping & 0,25 \\
redemption\_activation & 0,10 \\
product\_quality & 0,10 \\
customer\_service & 0,10 \\
refund\_return & 0,15 \\
usability & 0,10 \\
value\_price & 0,15 \\
fraud\_scam & 0,50 \\
other & 0,30 \\
\bottomrule
\end{tabular}
\end{table}

\section{Kết quả cho gói nâng cấp giáo trình}
\subsection{Baseline cổ điển cho các thuật toán cốt lõi}
Bảng \ref{tab:syllabus_baselines_vi} tổng hợp các baseline sentiment trong gói không thêm dữ liệu.

\begin{table}[H]
\centering
\caption{Baseline benchmark giáo trình trên test}
\label{tab:syllabus_baselines_vi}
\begin{tabular}{lccc}
\toprule
Mô hình & Recall\_0 & Precision\_0 & F2\_0 \\
\midrule
logreg\_l2 & 0,913 & 0,703 & 0,862 \\
multinomial\_nb & 0,913 & 0,633 & 0,839 \\
complement\_nb & 0,952 & 0,510 & 0,811 \\
perceptron & 0,828 & 0,684 & 0,795 \\
linear\_svm & 0,847 & 0,789 & 0,835 \\
sgd\_log\_loss & 0,921 & 0,616 & 0,838 \\
svd\_semantic\_logreg & 0,920 & 0,501 & 0,788 \\
ffnn\_mlp\_svd & 0,730 & 0,799 & 0,742 \\
\bottomrule
\end{tabular}
\end{table}

\subsection{Hành vi selective của các baseline bổ sung}
\begin{table}[H]
\centering
\caption{Chỉ số selective cho các baseline xác suất (test)}
\label{tab:syllabus_selective_vi}
\begin{tabular}{lcccc}
\toprule
Mô hình & Coverage & Sel. precision\_0 & Sel. recall\_0 & Sel. F2\_0 \\
\midrule
logreg\_l2\_selective & 0,887 & 0,763 & 0,946 & 0,903 \\
multinomial\_nb\_selective & 0,903 & 0,654 & 0,928 & 0,857 \\
complement\_nb\_selective & 0,900 & 0,545 & 0,963 & 0,835 \\
sgd\_log\_loss\_selective & 0,863 & 0,696 & 0,959 & 0,892 \\
\bottomrule
\end{tabular}
\end{table}

\subsection{Kết quả N-Gram Language Model}
\begin{table}[H]
\centering
\caption{Perplexity của n-gram LM (add-k smoothing)}
\label{tab:ngram_ppl_vi}
\begin{tabular}{lcc}
\toprule
Mô hình & Perplexity validation & Perplexity test \\
\midrule
Unigram add-k & 376,51 & 378,02 \\
Bigram add-k & 358,59 & 363,08 \\
\bottomrule
\end{tabular}
\end{table}

\subsection{Ma trận bao phủ học phần cập nhật}
Bảng \ref{tab:course_fit_vi} thể hiện mức bao phủ chủ đề học phần theo artifact hiện có. Điểm trung bình đạt \textbf{94,3\%}.

\begin{table}[H]
\centering
\caption{Ma trận mức độ bao phủ học phần từ artifact hiện trạng}
\label{tab:course_fit_vi}
\begin{tabular}{lc}
\toprule
Chủ đề & Bao phủ (\%) \\
\midrule
Probability & 100 \\
Regular Expressions & 50 \\
Evaluation Measures & 100 \\
N-Gram Language Models & 100 \\
Naive Bayes & 100 \\
Perceptron & 100 \\
Logistic Regression & 100 \\
Feed-forward Neural Networks & 100 \\
Pytorch Basic & 100 \\
Vector Semantics and Embeddings & 80 \\
Recurrent Neural Networks & 100 \\
Transformers and LLMs & 100 \\
Masked Language Models & 100 \\
Applications of LLMs & 90 \\
\bottomrule
\end{tabular}
\end{table}

\section{Thảo luận Classic so với Transformer}
Ở bài toán sentiment, transformer vượt trội về accuracy tổng thể và precision\_0 (0,799 so với 0,677), trong khi classic V6 tốt hơn về recall\_0 (0,932 so với 0,853) và F2\_0 (0,867 so với 0,842). Khi dùng cùng dải ngưỡng 0,40/0,60, transformer có selective precision cao hơn nhưng selective recall thấp hơn classic.

Với issue extraction, mô hình classic đa nhãn hiện vẫn mạnh hơn transformer issue (micro-F1 0,792 so với 0,646 theo file so sánh hybrid). Hybrid route cải thiện so với transformer thuần nhưng vẫn chưa vượt classic.

\section{Phân tích lỗi và failure modes}
Gói đánh giá nghiêm ngặt ghi nhận 1.454 mẫu sentiment dự đoán sai (194 false negative lớp 0 và 1.260 false positive lớp 0). Với false positive, nhóm lỗi lớn nhất là negation pattern (49,8\%), tiếp theo là ngữ cảnh khác (32,4\%) và contrast pattern (10,4\%). Với false negative, nhóm lớn nhất là ngữ cảnh khác (52,1\%), short text (23,2\%) và negation pattern (14,4\%).

Kiểm định cặp mô hình cho thấy V6 có chênh lệch dương so với một số baseline (ví dụ recall\_0 +0,063 so với multinomial NB, +0,085 so với LinearSVM), nhưng p-value vẫn lớn hơn 0,05 ở lần chạy này; do đó xem như xu hướng tích cực, chưa kết luận ý nghĩa thống kê mạnh.

\section{Giá trị thực tiễn cho học phần NLP}
Dự án hiện có chuỗi artifact tái lập đầy đủ: dữ liệu đã xử lý, quy trình theo bước, model lưu, phân tích ngưỡng, chính sách uncertain, extension transformer, và gói đánh giá nghiêm ngặt. Điều này đáp ứng tốt cả mục tiêu triển khai kỹ thuật lẫn yêu cầu bảo vệ học thuật.

Gói nâng cấp giáo trình cải thiện rõ mức bao phủ chủ đề mà không cần mở rộng dữ liệu, phù hợp với hướng đánh giá trọng tâm phương pháp.

\section{Hạn chế và hướng phát triển}
Các hạn chế chính gồm mất cân bằng lớp mạnh, nhãn \texttt{other} chiếm ưu thế trong issue, và hiệu năng transformer đa nhãn còn thấp với cấu hình fine-tune hiện tại trên tập nhỏ. Taxonomy lỗi cũng cho thấy hệ thống vẫn nhạy với negation và câu tương phản phức hợp.

Hướng tiếp theo nên tập trung vào ensemble có calibration, tăng cường ngữ cảnh không thêm nhãn mới, và các mô hình chuỗi gọn nhẹ cho issue extraction nhưng vẫn giữ khả năng diễn giải.

\section{Kết luận}
Bản báo cáo hiện tại tổng hợp cấu trúc và số liệu theo một mạch phân tích nhất quán. Mô hình production chính cho sentiment vẫn là classic V6 vì đáp ứng tiêu chí recall/F2 cho lớp negative theo mục tiêu dự án. Transformer sentiment có lợi thế rõ ở precision và accuracy tổng thể, phù hợp vai trò bổ trợ. Với issue extraction, OVR classic có calibration vẫn là lựa chọn mạnh nhất trong các thực nghiệm hiện tại. Gói nâng cấp không thêm dữ liệu đã nâng đáng kể tính học thuật, tính tái lập và độ phủ chủ đề của dự án.

\clearpage
\section*{Tài liệu tham khảo}
\addcontentsline{toc}{section}{Tài liệu tham khảo}

\begin{thebibliography}{99}

\bibitem{jurafsky2026}
Dan Jurafsky and James H. Martin.
\textit{Speech and Language Processing} (3rd ed. draft).
2026.
\newline\texttt{https://web.stanford.edu/~jurafsky/slp3/}

\bibitem{turney2002}
Peter D. Turney.
Thumbs Up or Thumbs Down? Semantic Orientation Applied to Unsupervised Classification of Reviews.
In \textit{Proceedings of ACL}, 2002.
\newline\texttt{https://aclanthology.org/P02-1053/}

\bibitem{wang2012}
Sida Wang and Christopher D. Manning.
Baselines and Bigrams: Simple, Good Sentiment and Topic Classification.
In \textit{Proceedings of ACL}, 2012.
\newline\texttt{https://aclanthology.org/P12-2018/}

\bibitem{tsoumakas2011}
Grigorios Tsoumakas, Ioannis Katakis, and Ioannis Vlahavas.
Random k-Labelsets for multilabel classification.
\textit{IEEE Transactions on Knowledge and Data Engineering}, 23(7):1079--1089, 2011.
\newline\texttt{https://doi.org/10.1109/TKDE.2010.164}

\bibitem{zhang2014}
Min-Ling Zhang and Zhi-Hua Zhou.
A review on multi-label learning algorithms.
\textit{IEEE Transactions on Knowledge and Data Engineering}, 26(8):1819--1837, 2014.
\newline\texttt{https://doi.org/10.1109/TKDE.2013.39}

\bibitem{vaswani2017}
Ashish Vaswani, Noam Shazeer, Niki Parmar, Jakob Uszkoreit, Llion Jones, Aidan N. Gomez, Lukasz Kaiser, and Illia Polosukhin.
Attention Is All You Need.
In \textit{Advances in Neural Information Processing Systems 30}, 2017.
\newline\texttt{https://arxiv.org/abs/1706.03762}

\bibitem{devlin2019}
Jacob Devlin, Ming-Wei Chang, Kenton Lee, and Kristina Toutanova.
BERT: Pre-training of Deep Bidirectional Transformers for Language Understanding.
In \textit{Proceedings of NAACL-HLT}, 2019.
\newline\texttt{https://aclanthology.org/N19-1423/}

\bibitem{saito2015}
Takaya Saito and Marc Rehmsmeier.
The precision-recall plot is more informative than the ROC plot when evaluating binary classifiers on imbalanced datasets.
\textit{PLOS ONE}, 10(3):e0118432, 2015.
\newline\texttt{https://doi.org/10.1371/journal.pone.0118432}

\bibitem{geifman2017}
Yonatan Geifman and Ran El-Yaniv.
Selective Classification for Deep Neural Networks.
\textit{arXiv preprint arXiv:1705.08500}, 2017.
\newline\texttt{https://arxiv.org/abs/1705.08500}

\end{thebibliography}

\clearpage
\section*{Phụ lục A: Hình sentiment}
\addcontentsline{toc}{section}{Phụ lục A: Hình sentiment}

\begin{figure}[H]
\centering
\includegraphics[width=0.82\textwidth]{\detokenize{results/dm2_steps/01_rating_distribution.png}}
\caption{Phân bố rating dữ liệu nguồn}
\end{figure}

\begin{figure}[H]
\centering
\includegraphics[width=0.82\textwidth]{\detokenize{results/dm2_steps/06_f1_vs_k.png}}
\caption{F1 validation theo Chi-square top K}
\end{figure}

\begin{figure}[H]
\centering
\includegraphics[width=0.82\textwidth]{\detokenize{results/dm2_steps/06_recall0_vs_k.png}}
\caption{Recall lớp negative theo top K}
\end{figure}

\begin{figure}[H]
\centering
\includegraphics[width=0.82\textwidth]{\detokenize{results/dm2_steps/06b_recall0_by_variant.png}}
\caption{Recall lớp negative theo biến thể cổ điển}
\end{figure}

\begin{figure}[H]
\centering
\includegraphics[width=0.82\textwidth]{\detokenize{results/dm2_steps/06b_precision0_by_variant.png}}
\caption{Precision lớp negative theo biến thể cổ điển}
\end{figure}

\begin{figure}[H]
\centering
\includegraphics[width=0.82\textwidth]{\detokenize{results/dm2_steps/09_prob_hist_test_strong.png}}
\caption{Histogram xác suất trên test strong-label}
\end{figure}

\begin{figure}[H]
\centering
\includegraphics[width=0.82\textwidth]{\detokenize{results/dm2_steps/10_threshold_tradeoff.png}}
\caption{Trade-off ngưỡng của mô hình cổ điển}
\end{figure}

\begin{figure}[H]
\centering
\includegraphics[width=0.82\textwidth]{\detokenize{results/nlp_ext/nlp_confusion_matrix.png}}
\caption{Confusion matrix của transformer sentiment}
\end{figure}

\begin{figure}[H]
\centering
\includegraphics[width=0.82\textwidth]{\detokenize{results/nlp_ext/nlp_prob_hist_test.png}}
\caption{Histogram xác suất của transformer sentiment}
\end{figure}

\begin{figure}[H]
\centering
\includegraphics[width=0.82\textwidth]{\detokenize{results/nlp_ext/nlp_threshold_tradeoff.png}}
\caption{Trade-off ngưỡng của transformer sentiment}
\end{figure}

\clearpage
\section*{Phụ lục B: Hình issue extraction}
\addcontentsline{toc}{section}{Phụ lục B: Hình issue extraction}

\begin{figure}[H]
\centering
\includegraphics[width=0.82\textwidth]{\detokenize{results/issue_steps/02_label_frequency.png}}
\caption{Tần suất nhãn issue}
\end{figure}

\begin{figure}[H]
\centering
\includegraphics[width=0.82\textwidth]{\detokenize{results/issue_steps/02_per_label_f1.png}}
\caption{F1 theo nhãn của OVR Logistic}
\end{figure}

\begin{figure}[H]
\centering
\includegraphics[width=0.82\textwidth]{\detokenize{results/issue_steps_char_demo/02_per_label_f1.png}}
\caption{F1 theo nhãn của biến thể word+char}
\end{figure}

\begin{figure}[H]
\centering
\includegraphics[width=0.82\textwidth]{\detokenize{results/nlp_ext/issue_transformer/nlp_issue_transformer_per_label_f1.png}}
\caption{F1 theo nhãn của transformer issue baseline}
\end{figure}

\clearpage
\section*{Phụ lục C: Hình gói nâng cấp giáo trình}
\addcontentsline{toc}{section}{Phụ lục C: Hình gói nâng cấp giáo trình}

\begin{figure}[H]
\centering
\includegraphics[width=0.82\textwidth]{\detokenize{results/nlp_ext/syllabus_upgrade/nlp_syllabus_bench_recall0.png}}
\caption{Recall\_0 test của các baseline giáo trình}
\end{figure}

\begin{figure}[H]
\centering
\includegraphics[width=0.82\textwidth]{\detokenize{results/nlp_ext/syllabus_upgrade/nlp_syllabus_bench_f2.png}}
\caption{F2\_0 test của các baseline giáo trình}
\end{figure}

\end{document}
